\section{Conclusion}

We added the output-cotangent second moment to a Muon-type optimizer and
measured what it does, across \nDistinctRuns{} full-budget training
trajectories and 28 learning-rate screening trials on one pinned recipe.

Three things are true of the modification at the recipe's learning rate. It
helps, by \mlpdiagselmag{} nats of held-out loss for the best geometry, in every
training window and on all three seeds we ran. It helps only as a complement to
input-side whitening --- substituting output factors for input factors is worse
than either alone, and degrading the input side is the most expensive mistake in
our arm family. And its value is governed by how current the statistic is
rather than how rich it is: exact diagonal factors recomputed every update beat
dense spectral factors that cannot be refreshed faster than once every four,
while deleting the same factors is worse than diagonalizing them.

Two things are true of the comparison. A matched null that performs identical
work with identity factors recovers a fifth to a third of the apparent gain, so
the attributable geometry effect is smaller than the headline, and the first arm
this project favoured was mostly execution path. And the baseline is not tuned:
its own learning rate is worth an order of magnitude more than the modification
is at a fixed rate, short-schedule tuning does not transfer to the full budget,
and at transferred rates two of three geometries lose to the baseline by more
than they ever won by. No arm reaches any loss target in less wall-clock time
than the baseline.

We think the direction remains worth pursuing, for a specific reason: the
geometry effect that survives the matched null is consistent, window-stable, and
seed-replicated, and the cost objection is attackable in ways the quality
objection would not be --- the expensive part of the output side turned out to be
the part worth discarding. But it is not yet a result, and the gap between
``promising'' and ``established'' here is not many more arms. It is one grafted
run to separate geometry from step size, one null at the selected rate, a
full-budget learning-rate neighborhood, one freshness ablation, and a set of
mechanistic probes that need no training at all --- nine runs, against the 54
configurations the design invites. We have specified them in
Section~\ref{sec:program} with decision rules fixed in advance, because the
effects at stake are close enough to this study's own control artifacts that the
order of operations determines the answer.

The transferable lesson is narrower than the method and more useful. Adding
machinery to an optimizer changes its numerics, its communication and its
scheduling, and those changes are not loss-neutral; a control that runs the
machinery and discards its output costs one run and told us that a fifth to a
third of our first effect was not what we thought. Comparing against an
untuned baseline is the older trap, and it cost us more: a fixed-rate
comparison was measuring a quantity an order of magnitude smaller than the one
a practitioner would care about. Neither discipline is expensive. Both were
decisive here.
